\subsubsection{比较强化学习训练配置与追加训练效果}\label{sec:q3-11}

首版选定的 PPO 模型包含 62882 个参数，以 60 维候选特征和 12 维全局上下文输入策略与价值网络。
选定模型在四个训练阶段依次采样 16384、12288、12288、16384 条训练轨迹，合计 57344 条；
四阶段均未使用行为克隆。
另有18项训练试验、53次开发评估用于比较候选配置，其总采样量与选定模型的训练量分别记录。

在相同基础模型、训练种子和新增 16384 条轨迹下，GAE λ=0.95 模型相较 λ=1 控制组节省 62.38 s，区间为 {[}19.88,104.96{]} s；
相较未追加训练的基础模型则节省 48.20 s，区间为 {[}−4.73,100.61{]} s，仍跨零。
该对照同时改变优势估计与价值网络的 λ-return 目标，反映这一训练配置的综合影响。
强注意力模型相较基础模型反而平均多用 25.98 s，扩展长轴探点后的追加训练、加入覆盖负债特征也未形成稳定收益。
它们属于开发集上的机制探索，说明更大的网络、更多动作候选和更久训练并不自动带来改善；
尚不能由此判定这些方法族无效。

后续纯 CPU 训练扩大了并行采样，批量也由 32 增为 128，并从同一基础模型进行不同种子的继续训练。
因此，现有记录不能将变化单独归因于``从 8 核增加到更多 CPU 核''；
核心数、采样预算、更新批量和优化过程同时变化，且逻辑 CPU 槽不等于物理核心。

将后续继续训练模型记为RL-2，已确认的状态搜索尾段精化版记为SS-2，并在同机、同引擎、同48个开发场景配对，RL 均时为 3110.20 s，SS-2为 3101.98 s。
RL 平均慢 8.22 s，按``状态用时减 RL 用时''计算的区间为 {[}−58.87,41.21{]} s，尚不能判断平均性能具有可靠差异。
RL 每局移动省 29.86 s，但测量多 7.75 次、耗时多 38.75 s，抵消路线收益。
该结果与首版独立测试属于不同模型、不同场景；
它提示进一步优化应关注信息获取与行程的联合代价，现有结果尚不足以判定两个方法族的最终优劣。

\subsubsection{分析强化学习扫描行为与状态搜索官方演练结果}\label{sec:q3-12}

后期RL-2的48条既有轨迹中，48/48 局均先在原点连续测量全部 20 个频道。
该模型所属的自由扫描版本已取消初始全扫硬规则，每次可选择覆盖点---频道对，出现正反馈后还可选择定位、清除或中断扫描；
因此，这一行为确由已训练策略选择。
不过，七个覆盖点、几何候选和可靠退出机制仍是预设的。
该现象支持在当前分布和候选空间内采用原点集中扫描的经验合理性，不能证明七点布局或状态压缩全局最优。
更早版本的原点扫描由程序固定，因此本项行为分析仅采用解除限制后的模型轨迹。

最后，以冻结的相对无信号组合版进行的既有 50 局官方演练均已完成，全清 657 个源；
本批预定跑满50局。
平均整局计费用时 2976.87 s，总时间除总源数为 226.55 s/源。
这批数据可作为真实接口兼容性与执行效果的补充，但未让各方案在同一官方场景配对，不能据此声称优于 RL 或将其均时与合成测试直接排名。
本文统一报告可直接核验的实际耗时；
各历史报告的下界口径在辅助证据索引中保留，不用于跨批次排名。

